\chapter{Свет}\label{ch:svet}

\textbf{Что вы узнете:} как менять цвет светодиода и почему цвет нужно «удерживать».

\section{Команда rgb}

На каждом дроне есть цветной светодиод. Цвет меняет команда \kod{rgb}. Вот шоу, в котором дрон по очереди горит красным, зелёным и синим:

\solfile{l03_light.py}

Цвет можно задать \emph{именем}. Всего восемь имён, у каждого есть номер:

\begin{center}
\begin{tabular}{@{}cccccccc@{}}
\toprule
0 & 1 & 2 & 3 & 4 & 5 & 6 & 7 \\
\kod{black} & \kod{red} & \kod{green} & \kod{yellow} & \kod{blue} & \kod{purple} & \kod{cyan} & \kod{white} \\
\bottomrule
\end{tabular}
\end{center}

Цвет \kod{black} выключает светодиод. Вместо имени можно написать номер: \kod{d.rgb(4)} — это синий.

\section{rgb занимает ноль секунд}

Это главное, что нужно запомнить про свет. Команда \kod{rgb} не длится: цвет меняется в этот момент и держится до следующего \kod{rgb}. Чтобы цвет можно было увидеть, после него нужен шаг, у которого есть длительность. В примере это \kod{d.set_z(100, 2)}: дрон остаётся на высоте 100\cm две секунды и светится.

\begin{vnimanie}
Если написать два \kod{rgb} подряд, первый цвет никто не увидит: между ними прошло ноль секунд.
\end{vnimanie}

\section{Свой цвет}

Кроме имён, цвет можно задать тремя числами от 0 до 255: яркость красного, зелёного и синего. Для этого есть \kod{d.rgb_value(255, 128, 0)} или \kod{d.rgb((255, 128, 0))} (это то же самое). Эта смесь даёт оранжевый.

\begin{zametka}
На живых дронах проверены только восемь именованных цветов. Свой цвет в симуляторе работает, но на дроне может выглядеть иначе.
\end{zametka}

Если написать неизвестное имя, например \kod{d.rgb('pink')}, будет ошибка: «Неизвестный цвет 'pink'. Доступны: black, red, green, yellow, blue, purple, cyan, white».

\begin{probuy}
  \item Светофор: красный 3 секунды, жёлтый 1 секунда, зелёный 3 секунды. Дрон висит на высоте 60\cm. Сколько секунд длится шоу?
  \item Радуга: покажите по очереди цвета с номерами от 1 до 7, по секунде каждый. Используйте номера, а не имена.
  \item Сделайте оранжевый свет через \kod{rgb_value(255, 128, 0)} на 3 секунды.
  \item Что увидит зритель в таком коде? Почему?\\
  \kod{d.rgb('red')}\quad\kod{d.rgb('green')}\quad\kod{d.set_z(100, 2)}
\end{probuy}
\otvety{3}
